\section{Threat Model and Security Requirements}
\label{sec:threats}

This section formalizes the threat model for multi-tenant LLM deployments, in which deployment and ownership are separated across administrative boundaries: it defines the system model, characterizes the adversary classes, and states the security requirements together with the attack vectors they must withstand.

\subsection{System Model}

The architecture is defined over a \emph{two-account boundary}. The \emph{Vendor Account} is trusted and controlled by the Software-as-a-Service (SaaS) vendor; it hosts the prompts, the cryptographic keys, the container registry, and the monitoring infrastructure. The \emph{Customer Account} is semi-trusted; it hosts the deployed agent, an Application Programming Interface (API) gateway, and an event-forwarding component. The boundary maps onto Amazon Web Services (AWS) account, Azure subscription or tenant, and Google Cloud Platform (GCP) project boundaries, so the model does not depend on a particular provider.

\subsection{Adversary Classes}

Three adversaries with increasing capabilities are considered. The \textbf{Customer Administrator} ($\mathcal{A}_1$) holds full privileges over the customer account: it can inspect deployed resources and container images, read environment variables, and intercept traffic. Its motivation ranges from curiosity to competitive intelligence or the replication of the agent without the vendor. The \textbf{External Attacker} ($\mathcal{A}_2$) reaches only the public API of the agent and relies on prompt injection, jailbreaking, and social engineering through the model's input--output interface. The \textbf{Malicious Insider} ($\mathcal{A}_3$) is a member of either organization with partial access, for instance to development environments, CI/CD pipelines, or logging systems, who attempts to exfiltrate prompt intellectual property.

\subsection{Security Requirements and Attack Mapping}

The architecture must provide three guarantees. \emph{Prompt confidentiality} requires that the plaintext prompt never be accessible inside the customer account, neither at rest nor in transit. \emph{Output integrity} requires that responses reflect the vendor's prompt and that any tampering be detectable. \emph{Forensic traceability} requires that every response carry an invisible per-request fingerprint, so that leaked content can be attributed to an account, a session, and a request; the fingerprint is an attribution mechanism built into the response path, and its robustness against Unicode normalization, deliberate removal, or collusion is not claimed here. Table~\ref{tab:threats} maps each attack vector to the adversaries able to execute it, to the guarantee it targets, and to the defense layer that counters it. The table uses the abbreviations Hash-based Message Authentication Code (HMAC), man-in-the-middle (MITM), and zero-width character (ZWC).

\begin{table}[t]
\caption{\label{tab:threats}Attack vectors, adversary capabilities, and defense layer mapping.}
\centering
\scriptsize
\setlength{\tabcolsep}{2pt}
\begin{tabular}{@{}lccl@{}}
\toprule
\textbf{Attack vector} & \textbf{Adv.} & \textbf{Goal} & \textbf{Defense} \\
\midrule
Console/API inspection       & $\mathcal{A}_1$                  & Conf. & L1: Cross-acct isolation \\
Container reverse eng.       & $\mathcal{A}_1$                  & Conf. & L1: Build-time secret \\
Direct gatekeeper call       & $\mathcal{A}_1$                  & Conf. & L1: HMAC + account binding \\
Request replay               & $\mathcal{A}_1, \mathcal{A}_2$   & Conf. & L1: Timestamp + nonce \\
MITM tampering               & $\mathcal{A}_1$                  & Integ. & L1: Body hash in HMAC \\
Prompt injection / jailbreak & $\mathcal{A}_2$                  & Conf. & L2: Risk scoring \\
Multi-turn extraction        & $\mathcal{A}_2$                  & Conf. & L2+L3: Score + filter \\
Output exfiltration          & $\mathcal{A}_2, \mathcal{A}_3$   & Conf. & L3: Canary + detection \\
Unauthorized redistribution  & $\mathcal{A}_3$                  & Trace. & L3: ZWC fingerprint \\
Encoding evasion             & $\mathcal{A}_2$                  & Conf. & L2: Encoding detectors \\
Typoglycemia attacks         & $\mathcal{A}_2$                  & Conf. & L2: Normalization \\
Monitoring evasion           & $\mathcal{A}_1$                  & All   & L4: Vendor-side logs \\
\bottomrule
\end{tabular}
\end{table}

\textbf{Trust assumptions.} The analysis rests on four assumptions: (1)~the vendor's Identity and Access Management (IAM) configuration is correct; (2)~all cross-account communication uses Transport Layer Security (TLS) 1.2 or later; (3)~the container runtime isolates process memory; and (4)~the LLM API does not retain prompts beyond the session. Assumption~(1) is weaker than the hardware attestation of a TEE, but it holds on every cloud provider without additional hardware. Assumption~(3) delimits the scope of the model: an adversary who attaches a debugger to the agent process, inspects its runtime memory, or otherwise compromises the process can read the decrypted prompt while it is in use, and no software-only layer of the architecture prevents this. Such an adversary is outside the threat model; the per-request fingerprint of Layer~3 makes the resulting disclosure attributable rather than preventable, and deployments that must resist it require hardware-backed isolation.
